\section{Normal Forms}

% ---------------------------------------------------------------------------
\begin{frame}{Why Normal Forms? -- The Problem}
  \begin{alertblock}{Starting Point: A ``Naive'' Table}
    A table that stores everything in one place leads to systematic problems.
  \end{alertblock}

  \vspace{0.4em}
  \scriptsize
  \begin{tabular}{llllllll}
    \toprule
    \textbf{ID} & \textbf{Title} & \textbf{Songs} & \textbf{Artist} & \textbf{Debut} & \textbf{Genre} & \textbf{Rel.} \\
    \midrule
    A100 & Midnights & Would've...; Karma & Taylor Swift & 2006 & Pop & 2022 \\
    A101 & 1989 (TV) & Shake It Off; Blank Space & Taylor Swift & 2006 & Pop & 2023 \\
    A102 & SOS & Anti-Hero; Die For You & SZA & 2012 & R\&B & 2022 \\
    \bottomrule
  \end{tabular}
  \normalsize

  \vspace{0.5em}
  \begin{block}{Spot the Problems}
    \begin{itemize}
      \item \textbf{Songs column} contains multiple values in one cell -- not atomic!
      \item \textbf{Debut year} 2006 appears twice for Taylor Swift -- redundancy!
      \item If Taylor Swift changes her name, we must update \emph{every} row.
    \end{itemize}
  \end{block}
\end{frame}

% ---------------------------------------------------------------------------
\begin{frame}{The Four Anomalies}
  \begin{block}{Update Anomaly}
    If Taylor Swift's debut year is stored in $n$ rows and we update only $k < n$,
    the database becomes \textbf{inconsistent} -- some rows say 2006, others say 2004.
  \end{block}

  \begin{block}{Insertion Anomaly}
    We \textbf{cannot record a new artist} unless they have already released an album.
    The schema forces us to have album data before artist data can be stored.
  \end{block}

  \begin{block}{Deletion Anomaly}
    If we delete SZA's only album (A102), we \textbf{lose all information about SZA}
    as well -- artist data exists only inside the album rows.
  \end{block}

  \begin{alertblock}{Root Cause}
    All anomalies stem from storing \textbf{multiple independent facts} in the same
    relation without separating their concerns.
  \end{alertblock}
\end{frame}

% ---------------------------------------------------------------------------
\begin{frame}{The Solution: Normal Forms}
  \begin{block}{What Are Normal Forms?}
    \textbf{Normal forms (NF)} are a series of progressive design rules.
    Each rule eliminates a specific class of redundancy or dependency problem.
  \end{block}

  \begin{block}{Key Properties}
    \begin{itemize}
      \item Normal forms are \textbf{cumulative}: a relation in $n$NF must also
            satisfy all lower normal forms -- there is no skipping.
      \item Each step decomposes relations into smaller, cleaner ones.
      \item The decomposition must be \textbf{lossless} (no information destroyed).
    \end{itemize}
  \end{block}

  \begin{columns}[t]
    \begin{column}{0.48\textwidth}
      \textbf{Normal Form Hierarchy}
      \begin{itemize}
        \item \textbf{1NF} -- atomic values
        \item \textbf{2NF} -- no partial FDs
        \item \textbf{3NF} -- no transitive FDs
        \item \textbf{BCNF} -- every determinant is a superkey
        \item \textbf{4NF} -- no multivalued FDs
        \item \textbf{5NF} -- no join dependencies
      \end{itemize}
    \end{column}
    \begin{column}{0.48\textwidth}
      \textbf{Two Key Algorithms}
      \begin{itemize}
        \item \textbf{Synthesis} $\to$ 3NF\\
              lossless + dependency-preserving
        \item \textbf{Decomposition} $\to$ BCNF\\
              lossless, may lose some FDs
      \end{itemize}
      \vspace{0.5em}
      \textbf{In practice:} Most production databases are designed to \textbf{3NF} or \textbf{BCNF}.
    \end{column}
  \end{columns}
\end{frame}
